\setcounter{section}{0}

\Needspace{5\baselineskip}
\hypertarget{ux5377ux79efux795eux7ecfux7f51ux7edccnn}{%
\section{卷积神经网络（CNN）}\label{ux5377ux79efux795eux7ecfux7f51ux7edccnn}}

\Needspace{5\baselineskip}
\hypertarget{ux4e00ux6838ux5fc3ux601dux60f3ux6a21ux4effux4ebaux7c7bux89c6ux89c9ux7684ux5c42ux6b21ux5316ux7279ux5f81ux63d0ux53d6}{%
\subsection{一、核心思想：模仿人类视觉的层次化特征提取}\label{ux4e00ux6838ux5fc3ux601dux60f3ux6a21ux4effux4ebaux7c7bux89c6ux89c9ux7684ux5c42ux6b21ux5316ux7279ux5f81ux63d0ux53d6}}

如识别一只猫：你首先注意到的是局部细节：边缘（胡须、耳朵轮廓）、斑点、颜色块。然后这些细节组合成更大的结构：眼睛、鼻子、爪子。最后，这些结构组合成整体概念：猫。

\Needspace{5\baselineskip}
\hypertarget{ux4e8cux6838ux5fc3ux7ec4ux4ef6}{%
\subsection{二、核心组件}\label{ux4e8cux6838ux5fc3ux7ec4ux4ef6}}

\Needspace{4\baselineskip}
\begin{enumerate}
\def\labelenumi{\arabic{enumi}.}
\tightlist
\item
  卷积层（Convolutional Layer）：特征探测器的滑动扫描
\end{enumerate}

\Needspace{8\baselineskip}
（1）通俗解释

想象你拿一个放大镜（称为``滤波器''或``卷积核''）在图像上从左到右、从上到下慢慢滑动。这个放大镜有特殊能力，能``点亮''它看到的特定模式（比如垂直线、水平线、特定颜色）。它在每个位置都计算一个数值，表示该位置与放大镜关注模式的匹配程度。所有位置计算完后，你就得到了一张新的``特征图''，这张图突出了原始图像中符合该放大镜关注模式的所有区域。CNN 通常有很多个不同的放大镜（不同的滤波器），每个关注不同的模式（边缘、纹理、形状等），从而生成多张不同的特征图。

本质上，卷积核的每一个数值都是一些输入和一些输出之间的共同权重。输出单元只与输入的一个局部区域相连（感受野），而非全连接。同一个滤波器在输入的不同位置使用相同的权重。这意味着无论特征出现在图像的哪个位置，都由同一个探测器检测。这是 CNN 具有平移不变性（位置不变性）的基础，也是减少参数的核心。

\Needspace{8\baselineskip}
（2）输入

每一次的输入是一个四维张量，按Pytorch的规则，第一个维度是batch\_size；第二个维度是input\_channels（输入通道数，如RGB图像第一个卷积层的输入通道数为3，代表3个颜色通道，后续卷积层的输入通道数则为前一层得到的特征数，卷积核的深度需与图像的通道数相同）；第三、四个维度是图像的高度和宽度input\_h和input\_w.

\Needspace{8\baselineskip}
（3）卷积核

又称为滤波器，是一个小的、可学习的权重矩阵。滤波器矩阵有3个维度，分别为高度kernel\_h，宽度kernel\_w，通道数input\_channels。一个卷积层通常有多个滤波器，其个数即为输出通道数output\_channels。每个滤波器负责检测一种特定的局部空间模式，对应一个输出通道。

\Needspace{8\baselineskip}
（4）卷积运算

输出矩阵是四维张量，第一个维度是batch\_size；第二个维度是output\_channels（卷积核个数）；第三、四个维度是output\_h和output\_w.

对于batch中的每一个样本：滤波器在输入上滑动（步长paddle），在每一个位置进行点积运算（Σ(卷积核的每一个元素*输入样本对应位置的元素)，然后再对各通道求和），然后可能还会加一个偏置，得到该样本该通道该位置的输出。

\Needspace{8\baselineskip}
（5）卷积运算的底层实现

\Needspace{5\baselineskip}
由于并行计算的效率远远高于用for循环，我们希望把卷积过程用矩阵乘法A\emph{B形式表征，A代表输入，B代表卷积核。由于输入是四维的、卷积核是三维的，我们会把部分维度合并展平。注意到矩阵乘法和点积的统一性，我们希望A的列数=B的行数=待点积的向量维度。容易发现这个待点积的向量维度是kernel\_h}kernel\_w*input\_channels，故构造A和B的行列如下：

A的行数=batch\_size（可以看成batch\_size个矩阵拼接）\emph{output\_h}output\_w（在多少个位置上进行卷积，即多少次卷积运算）

A的列数=B的行数=input\_channels\emph{kernel\_h}kernel\_w（卷积运算，按行、列、通道数求和）

B的列数=output\_channels（可以看成output\_channels个不同滤波器各自识别不同特征的矩阵拼接）

\Needspace{5\baselineskip}
可以看出B就是权重矩阵，但A不是输入矩阵，而是需要对输入矩阵进行操作：

第b个样本，卷积核左上角在{[}i,j{]}，从输入矩阵里裁input\_data{[}b,i:i+kernel\_h,j:j+kernel\_w,:{]}，把这个矩阵展平，然后把所有这样的矩阵拼接。当然，为了加速，我们事实上往往会用unfold函数，对于每个卷积核中心可达的位置（行），都有其各个通道上``卷积核覆盖区''的所有kernel\_h\emph{kernel\_w个元素（列）,使用unfold提取图像块：第一步在维度2(高度)上用一个kernel\_h}1的``核''滑动，把这些元素``收集到核中心''，这时核中心可达位置有output\_h*input\_w个，结果形状：{[}batch\_size, input\_channels, output\_h, input\_w, kernel\_h{]}；第二步在维度3(宽度)上操作第一步得到的矩阵，结果形状：{[}batch\_size, input\_channels, output\_h, output\_w, kernel\_h, kernel\_w{]}。

有时还会加偏置，每个卷积核（每个输出通道）只有一个偏置。

\Needspace{4\baselineskip}
\begin{enumerate}
\def\labelenumi{\arabic{enumi}.}
\setcounter{enumi}{1}
\tightlist
\item
  激活函数（Activation Function）：引入非线性
\end{enumerate}

卷积层计算出来的数值（特征匹配度）通常是线性的。但现实世界的数据和决策往往是高度非线性的。激活函数就像一个``开关''或``阈值处理器''，它决定一个特征是否足够强到被传递到下一层。它给网络增加了非线性的处理能力，使得网络能够学习复杂的映射关系。

\Needspace{4\baselineskip}
\begin{enumerate}
\def\labelenumi{\arabic{enumi}.}
\setcounter{enumi}{2}
\tightlist
\item
  池化层（又称汇聚层，Pooling Layer）：降采样与空间信息压缩
\end{enumerate}

池化层就像一个``信息浓缩器''。它在一个小区域（比如 2x2 方块）内，选取一个代表值（比如最大值或平均值）。输入特征图，指定池化窗口大小P\_h*P\_w和步长S（通常等于窗口大小，即不重叠），窗口在输入特征图上按步长S滑动，在每个窗口区域内计算一个代表值，构成输出特征图。这个操作是逐通道独立进行的，没有需要学习的参数。

\Needspace{8\baselineskip}
（1）池化操作分类

最大池化（Max Pooling）：取窗口内所有元素的最大值。

平均池化（Average Pooling）：取窗口内所有元素的平均值。

\Needspace{8\baselineskip}
（2）池化操作的好处

降维（抓大放小）：汇聚层把特征图分成一个个小区域（如2x2像素），然后只从这个区域里选一个``代表''（最大值或平均值）。这就像把一张高分辨率照片缩略成小图，虽然细节少了，但猫的轮廓、主要特征还在。这极大减少后续层需要处理的数据量，降低计算成本和内存占用，防止网络过于复杂。

增强鲁棒性（位置不敏感）：如果图片里的猫稍微挪动了一点位置（平移几个像素），卷积层输出的特征位置会跟着变。但汇聚层只关心这个小区域里的``最强音''（最大值）或``整体调性''（平均值），只要这个区域里最强的特征还是猫耳朵，汇聚后的输出就能保持稳定。这使网络对目标的微小平移、旋转、缩放（尺度）变化更不敏感（即具有平移不变性），提高了模型的泛化能力。

防止过拟合（模糊细节）：通过丢弃精确的位置信息和一些非关键细节，降低了网络对训练数据中特定噪声或无关细节的敏感性，使模型更关注是否存在某种特征，而不是特征具体在哪个精确像素点。

扩大感受野（看得更广）：随着网络加深，后续的卷积层作用在汇聚后的特征图上，相当于能``看到''原始输入图像中更大的区域（更大的感受野），有助于捕捉更高层次的语义信息（如整个猫头而不是单根胡须）

\Needspace{8\baselineskip}
（3）现代改进

可学习或平滑的汇聚规则：研究通过可学习参数或平滑加权，让汇聚不再固定为 max 或 mean。例如，SoftPool 根据激活值的指数权重对局部特征加权求和。

结合注意力机制：利用注意力机制动态调整汇聚区域的重要性权重，更智能地选择或融合信息。

空洞卷积：通过增大卷积核的感受野（在元素间插入空洞）来避免使用汇聚层进行下采样，常用于需要高分辨率输出的任务（如语义分割）。

\Needspace{4\baselineskip}
\begin{enumerate}
\def\labelenumi{\arabic{enumi}.}
\setcounter{enumi}{3}
\tightlist
\item
  全连接层（Fully Connected Layer）：综合决策
\end{enumerate}

在CNN的最后几层，经过多次卷积、激活、池化后得到的特征图（可能被展平成一维向量）会被送入一个或多个全连接层。这就像传统的神经网络（多层感知机MLP）。每个神经元都连接到前一层的所有神经元。它的作用是综合前面提取到的所有高级抽象特征，并根据这些特征做出最终的决策（比如这张图是``猫''还是``狗''的概率）。

\Needspace{4\baselineskip}
\begin{enumerate}
\def\labelenumi{\arabic{enumi}.}
\setcounter{enumi}{4}
\tightlist
\item
  其他重要概念
\end{enumerate}

（1）填充（Padding）：在图像边缘周围添加一圈像素（通常是0），这样卷积时边缘信息也能被作为卷积区域的中心，从而让这些信息和中间信息得到``同等待遇''的卷积处理，控制输出特征图的大小不变。

（2）步长（Stride）：滤波器在水平和垂直方向上每次移动的像素数。步长越大，滑动越快，输出特征图越小，感受野增加越快。

（3）批归一化（Batch Normalization）：在训练过程中，对每一层（通常在卷积层或全连接层之后，激活函数之前）的输入进行归一化处理（减均值除标准差），并引入可学习的缩放和平移参数。这能加速训练、缓解梯度消失/爆炸、允许使用更高的学习率、有一定的正则化效果。

\Needspace{5\baselineskip}
对于一个 mini-batch 的输入 \(B=\{x_1,\ldots,x_m\}\)，计算 mini-batch 均值和方差：

\[
\mu_B=\frac{1}{m}\sum_{i=1}^{m}x_i,\qquad
\sigma_B^2=\frac{1}{m}\sum_{i=1}^{m}(x_i-\mu_B)^2
\]

\Needspace{5\baselineskip}
归一化与缩放平移为：

\[
\hat{x}_i=\frac{x_i-\mu_B}{\sqrt{\sigma_B^2+\epsilon}},\qquad
y_i=\gamma\hat{x}_i+\beta
\]

其中 \(\epsilon\) 是小的常数，用于防止除零；\(\gamma\) 和 \(\beta\) 是可学习参数，用于恢复网络的表达能力。

（4）Dropout：在训练时，随机让一部分神经元暂时``失效''（输出置零）。这迫使网络不能过分依赖某些特定的神经元，增强了模型的泛化能力，减少过拟合。训练时，对每个神经元以概率p将其输出置零。测试时，所有神经元都激活，但输出值要乘以1-p（或训练时乘以1/(1-p)，测试时不变）以保持期望输出值不变。

\Needspace{4\baselineskip}
\begin{enumerate}
\def\labelenumi{\arabic{enumi}.}
\setcounter{enumi}{5}
\tightlist
\item
  多通道（Multi-Channel）
\end{enumerate}

事实上，输入数据或特征图不是一个单一的二维矩阵（\(H \times W\)），而是一个三维张量：\(H(\mathrm{Height}) \times W(\mathrm{Width}) \times C(\mathrm{Channels})\)。

对于输入图像，一个标准的RGB彩色图像的输入为H\emph{W}3矩阵，有3个通道：红色通道（R）、绿色通道（G）、蓝色通道（B）。每个通道都是一个二维网格（矩阵），网格中的每个值（像素）表示该位置对应颜色（R/G/B）的强度。灰度图像只有1个通道（亮度）。

对于在CNN的卷积层之后产生的特征图，其通道数表示不同类型的特征。例如，第一个卷积层输出的某个通道可能主要对``垂直边缘''敏感，另一个通道可能对``45度边缘''或``红色区域''敏感。随着网络加深，通道代表的特征越来越抽象（如``车轮''、``猫耳朵''的某种模式）。某个卷积层输出的特征图的尺寸为H\emph{W}C\_out，输出通道数C\_out是该层使用的卷积核的数量，也就是不同类型特征的个数。

\Needspace{5\baselineskip}
在位置(i, j)使用卷积核进行卷积的操作分为如下步骤：

（1）提取输入块：在输入X上，取出一个与卷积核大小相同的三维块X{[}i:i+K\_h, j:j+K\_w, :{]}。这个块的维度是(K\_h, K\_w, C\_in)。输入通道与卷积核深度必须相等（都为C\_in），这是多通道卷积能融合不同通道信息的基础。卷积核在通道维度上``扫描''并整合所有输入通道的信息。

（2）逐通道点积并求和：将卷积核与这个输入块进行逐元素相乘，得到一个同样大小的三维结果块(K\_h, K\_w, C\_in)。例如，在RGB输入中，一个检测``红色物体边缘''的卷积核，可能会在R通道赋予正权重，在G和B通道赋予负权重或零权重；在深层特征图中，一个检测``车轮''的卷积核，可能会组合来自不同底层特征通道（如``圆形''、``黑色区域''、``纹理''）的信息。。

（3）全局求和并加偏置：将上一步得到的三维结果块中的所有元素相加，加上一个标量偏置项b\_k（这个偏置是每个卷积核独有的）。这将卷积核在每个空间位置、每个输入通道上计算的相关性汇总成一个标量值。这个标量值代表了该卷积核所检测的特定特征（如某种边缘、纹理、颜色组合）在当前位置(i, j)的综合响应强度。

（4）得到输出值：这个(Sum + b\_k)的值，就是输出特征图O在位置(i, j)，对应第k个卷积核（即输出通道k）的值：O{[}i, j, k{]} = b\_k + sum\_\{di=0\}\^{}\{K\_h-1\} sum\_\{dj=0\}\^{}\{K\_w-1\} sum\_\{c=0\}\^{}\{C\_in-1\} ( X{[}i+di, j+dj, c{]} * K{[}di, dj, c, k{]} )。

\Needspace{5\baselineskip}
\hypertarget{ux4e09ux5178ux578bux67b6ux6784ux6a21ux5f0f}{%
\subsection{三、典型架构模式}\label{ux4e09ux5178ux578bux67b6ux6784ux6a21ux5f0f}}

\Needspace{5\baselineskip}
经典的 CNN 架构通常遵循以下模式：

\texttt{输入图像\ -\textgreater{}\ {[}{[}卷积层\ -\textgreater{}\ 激活函数{]}\ *\ N\ -\textgreater{}\ 池化层{]}\ *\ M\ -\textgreater{}\ 展平层\ -\textgreater{}\ {[}全连接层\ -\textgreater{}\ 激活函数{]}\ *\ K\ -\textgreater{}\ 输出层}

随着层数加深，空间分辨率逐渐降低（池化、大步长卷积），通道数逐渐增加（使用更多不同的滤波器），提取的特征从低级（边缘、角点、纹理）向高级（物体部件、完整物体、场景）演变。最后全连接层将提取到的高级特征进行整合，完成最终任务（分类、回归等）。

经典卷积神经网络示例：LeNet

\Needspace{5\baselineskip}
LeNet-5（经典版本）由 7 层组成，输入为 \(32\times 32\) 灰度图像，输出为 0-9 数字的概率分布：

\begin{longtable}[]{@{}
  >{\raggedright\arraybackslash}p{(\columnwidth - 8\tabcolsep) * \real{0.197}}
  >{\raggedright\arraybackslash}p{(\columnwidth - 8\tabcolsep) * \real{0.197}}
  >{\raggedright\arraybackslash}p{(\columnwidth - 8\tabcolsep) * \real{0.197}}
  >{\raggedright\arraybackslash}p{(\columnwidth - 8\tabcolsep) * \real{0.197}}
  >{\raggedright\arraybackslash}p{(\columnwidth - 8\tabcolsep) * \real{0.197}}@{}}
\toprule
层名称 & 操作 & 输入尺寸 & 输出尺寸 & 核心作用 \\
\midrule
\endhead
输入层 & 原始图像 & \(32\times 32\times 1\) & \(32\times 32\times 1\) & 接收灰度图像 \\
C1卷积层 & 6 个 \(5\times 5\) 卷积核 & \(32\times 32\times 1\) & \(28\times 28\times 6\) & 提取边缘/纹理特征 \\
S2池化层 & \(2\times 2\) 平均池化 & \(28\times 28\times 6\) & \(14\times 14\times 6\) & 降维，保留关键特征 \\
C3卷积层 & 16 个 \(5\times 5\) 卷积核 & \(14\times 14\times 6\) & \(10\times 10\times 16\) & 提取复杂特征（如数字局部结构） \\
S4池化层 & \(2\times 2\) 平均池化 & \(10\times 10\times 16\) & \(5\times 5\times 16\) & 进一步降维 \\
C5全连接层 & 展平 + 全连接 & \(5\times 5\times 16\rightarrow 400\) & 120 & 综合高级特征 \\
F6全连接层 & 全连接 & 120 & 84 & 非线性特征映射 \\
输出层 & 全连接 + Softmax & 84 & 10 & 生成分类概率 \\
\bottomrule
\end{longtable}

\Needspace{5\baselineskip}
\hypertarget{ux56dbcnnux7684ux4f18ux52bf}{%
\subsection{四、CNN的优势}\label{ux56dbcnnux7684ux4f18ux52bf}}

\begin{enumerate}
\def\labelenumi{\arabic{enumi}.}
\item
  参数共享：同一个滤波器在整个图像上滑动使用相同的权重，极大减少了需要学习的参数数量（与全连接网络相比）。这使得训练更高效，模型更小，需要的训练数据相对更少，且具有平移不变性。
\item
  局部连接：输出单元只与输入的一个局部区域连接，这符合图像数据的局部相关性特性（像素的语义通常由其邻近像素决定），也减少了参数。
\item
  层次化特征提取：自动从原始像素开始学习越来越复杂的特征（边缘 -\textgreater{} 纹理 -\textgreater{} 部件 -\textgreater{} 物体），无需人工设计特征。
\item
  平移不变性：由于参数共享和池化操作，网络对目标在图像中的位置变化具有一定的不敏感性。特征出现在哪里不重要，重要的是它出现了。
\item
  空间层次性：网络结构自然地通过卷积和池化操作建立了空间上的层次结构。
\end{enumerate}

\Needspace{5\baselineskip}
\hypertarget{ux4e94ux5377ux79efux795eux7ecfux7f51ux7edcux7684ux67b6ux6784ux53d8ux5316ux8d8bux52bf}{%
\subsection{五、卷积神经网络的架构变化趋势}\label{ux4e94ux5377ux79efux795eux7ecfux7f51ux7edcux7684ux67b6ux6784ux53d8ux5316ux8d8bux52bf}}

\begin{enumerate}
\def\labelenumi{\arabic{enumi}.}
\item
  减少池化层，用步长大于1的卷积代替降采样。
\item
  大量使用批量归一化层。
\item
  使用更深的网络（ResNet, DenseNet等通过跳跃连接解决深度网络训练难题）。
\item
  使用更小的滤波器（如3x3卷积堆叠代替大卷积核）。
\end{enumerate}

\FloatBarrier
\Needspace{5\baselineskip}
\hypertarget{ux53c2ux8003ux6587ux732e-11}{%
\subsection{参考文献}\label{ux53c2ux8003ux6587ux732e-11}}

\vspace{0.25em}
\begin{itemize}
\tightlist
\item
  LeCun, Y., Bottou, L., Bengio, Y., \& Haffner, P. (1998). \href{http://yann.lecun.com/exdb/publis/pdf/lecun-98.pdf}{Gradient-Based Learning Applied to Document Recognition}. Proceedings of the IEEE.
\item
  Krizhevsky, A., Sutskever, I., \& Hinton, G. E. (2012). \href{https://papers.nips.cc/paper/4824-imagenet-classification-with-deep-convolutional-neural-networks}{ImageNet Classification with Deep Convolutional Neural Networks}. NeurIPS 2012.
\item
  Stergiou, A., Poppe, R., \& Kalliatakis, G. (2021). \href{https://arxiv.org/abs/2101.00440}{Refining Activation Downsampling with SoftPool}. ICCV 2021.
\end{itemize}

\clearpage
